% Chapter 12: Doubly Robust Estimation
% A First Course in Causal Inference - Peng Ding
% Rewritten in Stein Style

\chapter{双重稳健估计}\label{ch:12}

\section{引言：一个估计量的双重保险}\label{sec:12-intro}

在第 \ref{ch:11} 章中，我们遇到了一个基本的张力：使用观测数据进行因果推断时，我们需要在两个识别公式之间做出选择。

第一个是\textbf{结果回归公式}（outcome regression）：
\[
\tau = E\{\mu_1(X)\} - E\{\mu_0(X)\},
\label{eq:12-outcome-reg}
\]
其中 $\mu_1(X) = E(Y \mid Z=1, X)$ 和 $\mu_0(X) = E(Y \mid Z=0, X)$ 分别是治疗组和对照组在给定协变量 $X$ 下的条件均值函数。这个公式的美妙之处在于：只要我们能正确建模 $\mu_1(X)$ 和 $\mu_0(X)$，就能一致地估计 $\tau$。但问题在于，如果模型错了，估计量就会产生偏倚。

第二个是\textbf{逆概率加权公式}（IPW）：
\[
\tau = E\left\{\frac{ZY}{e(X)}\right\} - E\left\{\frac{(1-Z)Y}{1-e(X)}\right\},
\label{eq:12-ipw}
\]
其中 $e(X) = \Pr(Z=1 \mid X)$ 是倾向得分。这个公式同样优雅：只要倾向得分模型正确，估计量就是一致的。但它的弱点同样明显——如果 $e(X)$ 模型错了，后果可能是灾难性的。

一个自然的问题应运而生：\textbf{能否构造一个估计量，只要两个模型中任意一个正确，就能保证一致性？}

这正是 James Robins 及其合作者（\citealp{Scharfstein:1999, BangRobins:2005}）提出的双重稳健估计量（doubly robust estimator）的核心思想。正如我们在下面将看到的，这个估计量巧妙地结合了结果回归和逆概率加权的优点。

\section{总体版本的 DR 估计量}\label{sec:12-population}

为了构造双重稳健估计量，我们同时引入两个工作模型（working models）：
\begin{itemize}
  \item \textbf{结果模型}：$\mu_1(X, \beta_1)$ 和 $\mu_0(X, \beta_0)$，分别建模 $E(Y(1) \mid X)$ 和 $E(Y(0) \mid X)$。当结果模型正确时，$\mu_1(X, \beta_1) = \mu_1(X)$。
  \item \textbf{倾向得分模型}：$e(X, \alpha)$，建模 $\Pr(Z=1 \mid X)$。当倾向得分模型正确时，$e(X, \alpha) = e(X)$。
\end{itemize}

在实践中，这两个模型都可能错误设定——这也是为什么我们称之为"工作模型"。

基于这两个工作模型，我们定义：
\begin{align}
\tilde\mu_1^{\mathrm{dr}} &= E\left[\frac{Z\{Y - \mu_1(X, \beta_1)\}}{e(X, \alpha)} + \mu_1(X, \beta_1)\right],
\label{eq:12-dr-mu1} \\
\tilde\mu_0^{\mathrm{dr}} &= E\left[\frac{(1-Z)\{Y - \mu_0(X, \beta_0)\}}{1-e(X, \alpha)} + \mu_0(X, \beta_0)\right].
\label{eq:12-dr-mu0}
\end{align}

这两个公式有另一种等价的书写方式，将它们重新整理为：
\begin{align}
\tilde\mu_1^{\mathrm{dr}} &= E\left[\frac{ZY}{e(X, \alpha)} - \frac{Z - e(X, \alpha)}{e(X, \alpha)} \mu_1(X, \beta_1)\right],
\label{eq:12-dr-mu1-alt} \\
\tilde\mu_0^{\mathrm{dr}} &= E\left[\frac{(1-Z)Y}{1-e(X, \alpha)} - \frac{e(X, \alpha) - Z}{1-e(X, \alpha)} \mu_0(X, \beta_0)\right].
\label{eq:12-dr-mu0-alt}
\end{align}

观察这两种形式，我们发现：
\begin{itemize}
  \item \cref{eq:12-dr-mu1} 和 \cref{eq:12-dr-mu0} 是对结果回归估计量的"增强"——在结果回归的基础上添加了逆概率加权项来修正残差。
  \item \cref{eq:12-dr-mu1-alt} 和 \cref{eq:12-dr-mu0-alt} 是对 IPW 估计量的"增强"——在逆概率加权的基础上添加了插补结果的修正项。
\end{itemize}
正因为这种双向增强的结构，双重稳健估计量也被称为\textbf{增强逆概率加权估计量}（augmented inverse propensity score weighting, AIPW）。

\section{双重稳健性定理}\label{sec:12-dr-theorem}

下面这个定理是双重稳健估计量的理论核心。

\begin{Theorem}[双重稳健性]\label{th:12-dr}
假设可忽略性 $Z \perp \{Y(1), Y(0)\} \mid X$ 和 overlap $0 < e(X) < 1$。
\begin{enumerate}
  \item 若 $e(X, \alpha) = e(X)$ 或 $\mu_1(X, \beta_1) = \mu_1(X)$ 之一成立，则 $\tilde\mu_1^{\mathrm{dr}} = E\{Y(1)\}$。
  \item 若 $e(X, \alpha) = e(X)$ 或 $\mu_0(X, \beta_0) = \mu_0(X)$ 之一成立，则 $\tilde\mu_0^{\mathrm{dr}} = E\{Y(0)\}$。
  \item 若 $e(X, \alpha) = e(X)$ 或 $\{\mu_1(X, \beta_1) = \mu_1(X), \, \mu_0(X, \beta_0) = \mu_0(X)\}$ 之一成立，则 $\tilde\mu_1^{\mathrm{dr}} - \tilde\mu_0^{\mathrm{dr}} = \tau$。
\end{enumerate}
\end{Theorem}

\textbf{为什么称为"双重"稳健？} 因为只要倾向得分模型和结果模型中任意一个正确，$\tilde\mu_1^{\mathrm{dr}} - \tilde\mu_0^{\mathrm{dr}}$ 就等于 $\tau$。这给了估计量两次"保险"的机会。\footnote{定理 \cref{th:12-dr} 的完整证明见附录 \cref{sec:appendix-12-dr-proof}。}

\section{样本版本的 DR 估计量}\label{sec:12-sample}

从总体版本 \cref{eq:12-dr-mu1} 和 \cref{eq:12-dr-mu0} 出发，我们可以用样本 analog 直接构造估计量。

\begin{Definition}[平均因果效应的双重稳健估计量]\label{def:12-dr}
基于数据 $(X_i, Z_i, Y_i)_{i=1}^n$，按以下步骤构造 $\tau$ 的双重稳健估计量：
\begin{enumerate}
  \item 拟合倾向得分：得到 $e(X_i, \hat\alpha)$；
  \item 拟合结果均值：得到 $\mu_1(X_i, \hat\beta_1)$ 和 $\mu_0(X_i, \hat\beta_0)$；
  \item 构造双重稳健估计量 $\hat\tau^{\mathrm{dr}} = \hat\mu_1^{\mathrm{dr}} - \hat\mu_0^{\mathrm{dr}}$，其中
\end{enumerate}
\[
\hat\mu_1^{\mathrm{dr}} = \frac{1}{n}\sum_{i=1}^n\left[\frac{Z_i\{Y_i - \mu_1(X_i, \hat\beta_1)\}}{e(X_i, \hat\alpha)} + \mu_1(X_i, \hat\beta_1)\right],
\]
\[
\hat\mu_0^{\mathrm{dr}} = \frac{1}{n}\sum_{i=1}^n\left[\frac{(1-Z_i)\{Y_i - \mu_0(X_i, \hat\beta_0)\}}{1-e(X_i, \hat\alpha)} + \mu_0(X_i, \hat\beta_0)\right].
\]
\end{Definition}

根据 \cref{def:12-dr}，我们还可以将双重稳健估计量写成更直观的形式：\footnote{推导见附录 \cref{sec:appendix-12-alt-forms}}
\[
\hat\tau^{\mathrm{dr}} = \hat\tau^{\mathrm{reg}} + \frac{1}{n}\sum_{i=1}^n\frac{Z_i\{Y_i - \mu_1(X_i, \hat\beta_1)\}}{e(X_i, \hat\alpha)} - \frac{1}{n}\sum_{i=1}^n\frac{(1-Z_i)\{Y_i - \mu_0(X_i, \hat\beta_0)\}}{1-e(X_i, \hat\alpha)}.
\label{eq:12-dr-reg-aug}
\]

这个形式揭示了双重稳健估计量的本质：它从结果回归估计量 $\hat\tau^{\mathrm{reg}}$ 出发，通过添加修正项来获得鲁棒性。类似地，利用 \cref{eq:12-dr-mu1-alt} 和 \cref{eq:12-dr-mu0-alt}，我们也可以写成：\footnote{推导见附录 \cref{sec:appendix-12-alt-forms}}
\[
\hat\tau^{\mathrm{dr}} = \hat\tau^{\mathrm{ipw}} - \frac{1}{n}\sum_{i=1}^n\frac{Z_i - e(X_i, \hat\alpha)}{e(X_i, \hat\alpha)}\mu_1(X_i, \hat\beta_1) + \frac{1}{n}\sum_{i=1}^n\frac{e(X_i, \hat\alpha) - Z_i}{1-e(X_i, \hat\alpha)}\mu_0(X_i, \hat\beta_0).
\label{eq:12-dr-ipw-aug}
\]

关于方差估计，\cite{Funk:2011} 建议使用非参数 bootstrap 来近似 $\hat\tau^{\mathrm{dr}}$ 的方差。

\section{更深层的直觉}\label{sec:12-intuition}

虽然 \cref{eq:12-dr-mu1} 和 \cref{eq:12-dr-mu0} 的形式看起来并不显然，但它们可以从两个不同的直观角度推导出来。这两个角度恰好对应了结果回归和 IPW 的各自弱点。

\subsection{从 IPW 出发：降低方差}\label{sec:12-ipw-to-dr}

IPW 估计量 $\mu_1 = E\left\{\frac{ZY}{e(X)}\right\}$ 完全忽略了 $Y$ 的结果模型。它的优势在于不需要任何结果模型假设就能保持一致性。然而，如果协变量能够预测结果，那么即使工作结果模型错误，基于残差的伪潜在结果通常比原始结果具有更小的方差。

给定一个可能错误设定的工作模型 $\mu_1(X, \beta_1)$，我们可以做一个简单的分解：\footnote{推导见附录 \cref{sec:appendix-12-ipw-decomp}}
\[
\mu_1 = E\{Y(1)\} = E\{Y(1) - \mu_1(X, \beta_1)\} + E\{\mu_1(X, \beta_1)\} = E\left\{\frac{Z\{Y - \mu_1(X, \beta_1)\}}{e(X)} + \mu_1(X, \beta_1)\right\}.
\]
使用工作模型来提高效率是调查采样中的传统思想。\cite{LittleAn:2004} 和 \cite{Lumley:2011} 都指出了它与双重稳健估计量的联系。

\subsection{从结果回归出发：降低偏倚}\label{sec:12-reg-to-dr}

另一种思路是从结果回归出发，改善其偏倚。结果回归估计量 $\tilde\mu_1 = E\{\mu_1(X, \beta_1)\}$ 可能不等于 $\mu_1$，因为结果模型可能是错误的。这个估计量的偏倚是 $E\{\mu_1(X,\beta_1) - Y(1)\}$，它可以用 IPW 估计量来估计：\footnote{推导见附录 \cref{sec:appendix-12-bias-correct}}
\[
B = E\left\{\frac{Z\{\mu_1(X, \beta_1) - Y\}}{e(X)}\right\},
\]
只要倾向得分模型正确。因此，去偏估计量 $\tilde\mu_1 - B$ 恰好等于 \cref{eq:12-dr-mu1}。这给出了构造双重稳健估计量的第二个直观视角。

\section{模拟研究}\label{sec:12-simulation}

为了评估双重稳健估计量在有限样本下的表现，我们考虑四种场景：\footnote{模拟细节和完整结果见附录 \cref{sec:appendix-12-simulation}。}
\begin{enumerate}
  \item 倾向得分模型和结果模型都正确；
  \item 倾向得分模型错误，但结果模型正确；
  \item 倾向得分模型正确，但结果模型错误；
  \item 倾向得分模型和结果模型都错误。
\end{enumerate}

\textbf{关键发现}：场景 1-3 验证了双重稳健性——只要其中一个模型正确，DR 估计量就几乎无偏。但在场景 4（两个模型都错误）中，DR 估计量可能表现出"双重脆弱性"——两个误差的乘积导致比单独使用任一模型更大的偏差。\footnote{模拟结果表格见附录 \cref{sec:appendix-12-simulation-table}。}

\section{进一步讨论：为何脆弱？}\label{sec:12-discussion}

回顾 \cref{th:12-dr} 的证明，双重稳健性的关键在于以下乘积结构：\footnote{详细推导见附录 \cref{sec:appendix-12-product-structure}}
\[
\tilde\mu_1^{\mathrm{dr}} - E\{Y(1)\} = E\left[\frac{e(X) - e(X, \alpha)}{e(X, \alpha)} \times \{\mu_1(X) - \mu_1(X, \beta_1)\}\right].
\]
只要 $e(X) = e(X, \alpha)$ 或 $\mu_1(X) = \mu_1(X, \beta_1)$ 之一成立，上式就为零。然而，当两个模型都错误时，两个误差的乘积可能放大偏差——这就是"双重脆弱性"的数学本质。模拟研究证实了这一点。

尽管 \cite{KangSchafer:2007} 基于模拟研究批评了双重稳健估计量，但自 \cite{Scharfstein:1999} 的开创性工作以来，双重稳健估计量一直是因果推断的标准策略。近年来，它在理论统计和计量经济学文献中以更响亮的名字"双机器学习"（double machine learning）重新兴起。\cite{Chernozhukov:2018} 的基本思想是用更灵活的机器学习工具替代传统的参数工作模型。

\section{本章小结}\label{sec:12-summary}

本章介绍了双重稳健估计量，主要内容如下：

\begin{enumerate}
  \item \textbf{核心思想}：结合结果回归和 IPW，只要其中一个模型正确就能保证一致性。
  \item \textbf{公式结构（\cref{eq:12-dr-mu1}--\cref{eq:12-dr-mu0}）}：增强逆概率加权估计量（AIPW）。
  \item \textbf{双重稳健性（\cref{th:12-dr}）}：只要倾向得分模型或结果模型之一正确，$\tilde\mu_1^{\mathrm{dr}} - \tilde\mu_0^{\mathrm{dr}} = \tau$。
  \item \textbf{两种直觉视角}：从 IPW 降方差（\cref{sec:12-ipw-to-dr}）和从结果回归降偏倚（\cref{sec:12-reg-to-dr}）。
  \item \textbf{双重脆弱性}：当两个模型都错误时，两个误差的乘积可能放大偏差。
\end{enumerate}

\section{习题}\label{sec:12-homework}

\begin{HomeworkProblem}[A sanity check]\label{exr:12-1}
Consider the case in which the covariate is discrete $X \in \{1, \ldots, K\}$ and the parameter of interest is $\tau$. Without imposing any model assumptions, the estimated propensity score $\hat{e}(X)$ equals $\hat{e}_{[k]} = \hat{\pr}(Z=1\mid X=k) $, the proportion of units receiving the treatment, and the estimated outcome means are the sample means of the outcomes $\hat{\bar{Y}}_{[k]}(1)  = \hat{E}(Y\mid Z=1, X=k) $ and $\hat{\bar{Y}}_{[k]}(0)  = \hat{E}(Y\mid Z=0, X=k) $ under treatment, within stratum $X=k\ (k=1, \ldots, K)$. Show that the stratified estimator, outcome regression estimator, HT estimator, Hajek estimator, and doubly robust estimator are all identical numerically.



\end{HomeworkProblem}

\begin{HomeworkProblem}[An alternative form of the doubly robust estimator for $\tau$]\label{exr:12-2}

Motivated by \eqref{eq::explain-ipw-efficiency-basic}, we have an alternative form of the doubly robust estimator for $\mu_1 = E\{ Y(1) \} $:
\[
\tilde{\mu}_1^\text{dr2} =
\frac{ E\left[   \frac{Z\{ Y-\mu_1(X,\beta_1) \}}{  e(X,\alpha) } \right] }{E\left[   \frac{Z }{  e(X,\alpha) } \right]}
 + E\{ \mu_1(X, \beta_1)  \} .
\]
Show that $\tilde{\mu}_1^\text{dr2} = \mu_1$ if either $e(X,\alpha)  = e(X)$ or $\mu_1(X, \beta_1) =  \mu_1(X)$. Give the analogous formula for estimating $\mu_0$. Give the sample analog of the doubly robust estimator for $\tau$ based on these formulas.

Remark:
This form of doubly robust estimator appeared in \citet{robins2007comment}.



\end{HomeworkProblem}

\begin{HomeworkProblem}[An upper bound of the bias of the doubly robust estimator]\label{exr:12-3}


Consider the population version of the doubly robust estimator $\tilde{\mu}_1^\textup{dr} $ for $E\{  Y(1) \} $.
Show that
\[
 |   \tilde{\mu}_1^\textup{dr} - E\{  Y(1) \}  |  \leq
\sqrt{
E\left[   \frac{ \{  e(X) - e(X,\alpha)\}^2  }{  e(X,\alpha)^2 }  \right] \times
E\left[  \left\{  \mu_1(X) -  \mu_1(X, \beta_1)   \right\}^2 \right] .
}
\]
Find the analogous upper bound for the bias of $\tilde{\mu}_0^\textup{dr} $ for $E\{  Y(0) \} $.


Remark: You may find Section \ref{sec::cauchy-schwarz-inequality} useful for the proof.


\end{HomeworkProblem}

\begin{HomeworkProblem}[Data analysis of Example \ref{eg::job-training-obs}]\label{exr:12-4}


Analyze the dataset \ri{cps1re74.csv} using the methods discussed so far.



\end{HomeworkProblem}

\begin{HomeworkProblem}[Analyzing a dataset from the Karolinska Institute]\label{exr:12-5}

 \citet{rubin2008objective} used the dataset \texttt{karolinska.txt} to illustrate the ideas of causal inference in observational studies.
 The dataset has 158 cardia cancer patients diagnosed between 1988 and 1995 in Central and Northern Sweden, 79 diagnosed at large volume hospitals, defined as treating more than ten patients with cardia cancer during that period, and 79 diagnosed at the remaining small volume hospitals. The treatment \texttt{z}  is the indicator of whether a patient was diagnosed at a large volume hospital. The outcome \texttt{y} is whether the patient survived longer than 1 year after the diagnosis. The covariates \texttt{x} contain information about age, whether a patient was from a rural area, and whether a patient was male.

\begin{lstlisting}
karolinska = read.table("karolinska.txt", header = TRUE)
z = karolinska$hvdiag
y = 1 - (karolinska$year.survival == 1)
x = as.matrix(karolinska[, c(3, 4, 5)])
\end{lstlisting}


Analyze the dataset using the methods discussed so far.



\end{HomeworkProblem}

\begin{HomeworkProblem}[Recommended reading]\label{exr:12-6}

 \citet{lunceford2004stratification} gave a review and comparison of many methods discussed in Chapters \ref{chapter::pscore-key} and \ref{chapter::doubly-robust}.



\end{HomeworkProblem}
